\subsection{A moving kernel forces stable loss one}
\label{subsec:Cartier-kernel-wall-inertia}

Work over a perfect field of characteristic two; geometric conclusions
may be checked after passing to its algebraic closure.  All Frobenius
pullbacks below are typed explicitly.  This subsection supplies a local
criterion.  Its last remark states the nonemptiness condition still
needed for a new higher-degree application.

\begin{lemma}[Maximal kernel variation excludes a further nilpotent step]
\label{lem:maximal-kernel-variation-stable-loss}
\status{Candidate; complete proof supplied.}
Let \(Y\) be smooth, let \(D\subset Y\) be an integral divisor, and
let \(\mathcal E\) be a rank-\(r\) vector bundle with a regular
linearized Frobenius map
\[
             F:F_Y^*\mathcal E\longrightarrow\mathcal E.
                                                        \tag{KWI1}
\]
Near the generic point of \(D\), suppose that \(\mathcal E\) is a
subbundle of a constant bundle \(V\otimes\mathcal O_Y\), and that
the restriction of \(F\) to \(D\) has rank \(r-1\).
Let \(\mathcal K\) be its kernel line in \(F_D^*\mathcal E_D\).
If the corresponding rational map
\[
             D\dashrightarrow\mathbf P(V^{(2)})
                                                        \tag{KWI2}
\]
has differential rank \(r-1\) at the generic point, then the stable
Frobenius rank on \(D\) is exactly \(r-1\).
\end{lemma}
\begin{proof}
Shrink the smooth locus of \(D\) so that its first image
\(\mathcal I\subset\mathcal E_D\) is a rank-\(r-1\) subbundle.
At the generic point, Frobenius on that image is the restriction
\[
             F_D^*\mathcal I\longrightarrow\mathcal I.
                                                        \tag{KWI3}
\]
Its kernel is \(\mathcal K\cap F_D^*\mathcal I\).  Consequently
(KWI3) fails to be an isomorphism precisely when
\(\mathcal K\subset F_D^*\mathcal I\).

Suppose this containment holds.  Choose a rational frame
\(w_1,\ldots,w_{r-1}\) of \(\mathcal I\), considered as vectors
in the constant ambient bundle.  Their Frobenius pullbacks have
coordinate vectors \(w_i^{[2]}\).  Every coordinate of every such
vector has zero K\"ahler differential over the perfect ground
field.  A generator of \(\mathcal K\) can therefore be written
\[
                   v=\sum_{i=1}^{r-1}a_iw_i^{[2]},
       \qquad dv=\sum_{i=1}^{r-1}(da_i)w_i^{[2]}.
                                                        \tag{KWI4}
\]
Modulo the line \(kv\), the differential of its projective class
has image in \(F_D^*\mathcal I/\mathcal K\), of dimension
\(r-2\).  This contradicts the stipulated rank \(r-1\).
The containment is thus impossible.  Map (KWI3) is an isomorphism
at the generic point, and every subsequent image has the same
dimension \(r-1\), as asserted.  The proof uses the actual squared
coordinates in the constant ambient bundle; it does not assume
that the varying image subspace itself is constant.
\end{proof}

\begin{proposition}[Full hyperplane inertia on the common Cartier open]
\label{prop:birational-Cartier-kernel-wall-inertia}
\status{Candidate; complete proof supplied.}
Let \(W\) have odd dimension \(N\ge3\), put \(W'=W^{(2)}\),
and let
\[
                  b:\Lambda^2W\longrightarrow W'
                                                        \tag{KWI5}
\]
be an alternating bilinear map.  For \(s\in W\), write
\(M_s(t)=b(s,t)\), and suppose that \(U\subset\mathbf P(W)\)
is a nonempty open with the following properties:
\begin{enumerate}
\item \(M_s\) has rank \(N-1\) on \(U\), with left kernel
      \(k\lambda_s\subset(W')^*\);
\item the form \(B_{\lambda_s}(u,v)=\lambda_s(b(u,v))\) has
      rank \(N-1\) there;
\item on coefficient charts the relative coherent Frobenius is
      the restriction
      \[
       F_s=M_s^t:(s^{[2]})^\perp\longrightarrow s^\perp
      \]
      of the transpose, on the rank-\(N-1\) Euler bundle.
\end{enumerate}
Set
\[
                  H_U=\{[s]\in U:\lambda_s(s^{[2]})=0\}.
                                                        \tag{KWI6}
\]
If \(H_U\ne\varnothing\), it is a smooth reduced divisor in
\(U\).  At the geometric generic point of each of its components,
the relative Frobenius has first and stable rank \(N-2\), and its
determinant has order one.  The geometric divisorial inertia of the
generic relative Lang local system contains the full group
\[
          \operatorname{Hom}(\mathbf F_2,\mathbf F_2^{N-2})
                                                        \tag{KWI7}
\]
fixing its regular \((N-2)\)-dimensional unit subspace and its
one-dimensional quotient.
\end{proposition}
\begin{proof}
All assertions are local on coefficient charts, where the
tautological line is trivialized and \(s\) is a regular nonzero
vector.  Alternation gives \(M_ss=0\); its stipulated rank makes
\(ks\) the full right kernel.  Let
\(\kappa([s])=[\lambda_s]\).  Differentiating
\(\lambda_sM_s=0\) gives
\[
          (d\lambda_s(t))M_s+\lambda_sM_t=0.
                                                        \tag{KWI8}
\]
The isomorphism \(W/ks\xrightarrow{M_s}\ker\lambda_s\)
identifies the derivative of \(\kappa\) with the alternating
form \(B_{\lambda_s}\) on \(W/ks\).  Its rank is \(N-1\),
so \(d\kappa\) is an isomorphism at every point of \(U\).
Equivalently, \(\kappa\) is etale there.  This argument does not
identify \(M_s\) itself with an alternating matrix.

Choose a local generator \(\lambda_s\).  Put
\(q(s)=\lambda_s(s^{[2]})\).  At a point where \(q=0\),
\[
                 dq(t)=(d\lambda_s(t))(s^{[2]}).
                                                        \tag{KWI9}
\]
Evaluation on \(s^{[2]}\) is a nonzero functional on
\((W')^*/k\lambda_s\), since \(s^{[2]}\ne0\) and
\(\lambda_s(s^{[2]})=0\).  The isomorphism \(d\kappa\)
therefore makes (KWI9) nonzero.  Rescaling the generator multiplies
this differential on \(q=0\) by a unit.  Thus (KWI6) is a smooth
reduced divisor wherever it is nonempty.

The full transpose induces an isomorphism
\((W')^*/k\lambda_s\simeq s^\perp\).  Restricting its source
to \((s^{[2]})^\perp\) gives an isomorphism off \(q=0\),
and has the one-dimensional kernel \(k\lambda_s\) on that
divisor.  In local frames its determinant is a unit times \(q\):
this is the determinant of the projection of a hyperplane along
the line \(k\lambda_s\), whose evaluation on the defining
covector \(s^{[2]}\) is exactly \(q\).  Hence the order is one,
not merely a positive order inferred from the rank.

On a component \(D\) of \(H_U\), the linearized kernel map in
the constant ambient \((W')^*\) is exactly
\(\kappa|_D\).  Its differential has rank \(N-2\), since
\(\kappa\) is etale.  On each coefficient chart the Euler
bundle is a subbundle of the constant \(W^*\)-bundle; its
Frobenius pullback embeds in the constant \((W')^*\)-bundle.
The preceding lemma, with \(r=N-1\), now gives stable rank
\(N-2\).  All hypotheses of
Lemma~\ref{lem:moving-kernel-full-root-inertia-candidate} hold:
simple determinant order, one kernel line, stable loss one, and
maximal kernel differential rank.  Its complete strict-henselian
Artin--Schreier proof gives (KWI7), including possible inseparable
residue extensions.  Thus the conclusion is the full hyperplane
radical, not an assumption that inertia is cyclic.
\end{proof}

\begin{remark}[The remaining nonemptiness and generation conditions]
\label{rem:Cartier-kernel-wall-inertia-scope}
\status{Candidate scope; no new unconditional monodromy degree.}
The common Cartier-incidence argument using the typed Moore slices
produces an open \(U\) satisfying the first two conditions above
for odd etale pullbacks of cubic principal polarizations.  The
universal hypersurface and Euler sequences give the third.  The
proposition still requires \(H_U\ne\varnothing\).

On the nine-dimensional trivial-character section space, the
degree-one symmetrizer makes \(\lambda_s\) a fixed invertible
linear transform of \(s\), and (KWI6) becomes the previously
proved smooth degree-one cubic.  The finitely many nontrivial
character multiplication and dual-contraction determinants are
nonzero polynomials on that space.  Their nonvanishing on the
ambient space alone does not prove that they are nonzero on this
cubic: one must exclude its being a factor of each relevant
determinant.  A valid exclusion gives \(H_U\ne\varnothing\)
by intersecting finitely many nonempty opens of that integral
cubic.  No such universal exclusion is asserted here.

Once nonemptiness holds, a separate global argument is still needed
to generate the entire monodromy image from full hyperplane
radicals, or to prove connectedness of the nonzero Lang cover.
Containment of one full radical and absence of invariant vectors
or covectors alone does not exclude the previously retained
parabolic countermodels.  In particular this subsection makes no
new fixed-surface full-image claim.  Its proofs are direct; no new
numerical test or primary-source assertion is used.  The complete
earlier moving-kernel proof has been reread for the stated local
application.
\end{remark}
